\documentclass[11pt,a4paper]{article}
\newcommand{\gradelabel}{Grade 12}
\newcommand{\leconhdr}{Architecture d'un objet connecté}
% Préambule commun aux plans de cours de robotique (Grades 7 à 12).
% Chaque plan de cours commence par \documentclass[11pt,a4paper]{article}, définit
% \gradelabel et \leconhdr (en-têtes de page), puis % Préambule commun aux plans de cours de robotique (Grades 7 à 12).
% Chaque plan de cours commence par \documentclass[11pt,a4paper]{article}, définit
% \gradelabel et \leconhdr (en-têtes de page), puis % Préambule commun aux plans de cours de robotique (Grades 7 à 12).
% Chaque plan de cours commence par \documentclass[11pt,a4paper]{article}, définit
% \gradelabel et \leconhdr (en-têtes de page), puis \input{../../preambule_robotique.tex}.
\usepackage[T1]{fontenc}
\usepackage[utf8]{inputenc}
\usepackage{lmodern}
\usepackage[french]{babel}
\usepackage[margin=2cm,headheight=14pt]{geometry}
\usepackage[most]{tcolorbox}
\usepackage{xcolor,tabularx,booktabs,array,enumitem,fancyhdr,amssymb,tikz}
\usetikzlibrary{arrows.meta,positioning}

\pagestyle{fancy}
\fancyhf{}
\lhead{\small Robotique -- \gradelabel{} -- Novembre 2026}
\rhead{\small \leconhdr}
\cfoot{\thepage}
\setlist{itemsep=1pt,topsep=2pt,leftmargin=*}
\renewcommand{\arraystretch}{1.35}
\newcolumntype{Y}{>{\raggedright\arraybackslash}X}
\newcommand{\cb}{$\square$\,}
\newcommand{\fl}{$\rightarrow$}
\newcommand{\lig}{\noindent\rule{\linewidth}{0.3pt}\\[5mm]}

\newtcolorbox{etape}[3]{enhanced,breakable,colback=#1!6,colframe=#1!70!black,
  coltitle=white,fonttitle=\bfseries,title={#2\hfill #3},boxrule=0.8pt,arc=2mm,
  before skip=8pt,after skip=8pt}
\newtcolorbox{astuce}[1][Différenciation et besoins particuliers]{breakable,
  colback=yellow!12,colframe=yellow!50!black,fonttitle=\bfseries\small,
  title={#1},boxrule=0.5pt,arc=1mm,before skip=6pt,after skip=6pt}
\newtcolorbox{code}[1][Code]{breakable,colback=black!4,colframe=black!60,
  fonttitle=\bfseries\small,title={#1},boxrule=0.5pt,arc=1mm,fontupper=\ttfamily\small,
  before skip=6pt,after skip=6pt}
% Blocs de code coloré (écrire le code sans accents : comme dans l'éditeur)
\newtcblisting{arduino}[1][Code Arduino]{listing only,breakable,colback=black!4,
  colframe=black!60,fonttitle=\bfseries\small,title={#1},boxrule=0.5pt,arc=1mm,
  before skip=6pt,after skip=6pt,
  listing options={language=C++,basicstyle=\ttfamily\small,columns=fullflexible,
  keepspaces=true,showstringspaces=false,commentstyle=\color{green!40!black}\itshape,
  keywordstyle=\color{blue!70!black}\bfseries}}
\newtcblisting{python}[1][Code Python]{listing only,breakable,colback=black!4,
  colframe=black!60,fonttitle=\bfseries\small,title={#1},boxrule=0.5pt,arc=1mm,
  before skip=6pt,after skip=6pt,
  listing options={language=Python,basicstyle=\ttfamily\small,columns=fullflexible,
  keepspaces=true,showstringspaces=false,commentstyle=\color{green!40!black}\itshape,
  keywordstyle=\color{blue!70!black}\bfseries}}
\newcommand{\entete}[2]{%
  \begin{tcolorbox}[colback=#1!10,colframe=#1!70!black,boxrule=0.8pt,arc=2mm]
  {\Large\bfseries #2}\hfill Nom : \rule{4cm}{0.3pt}\quad Date : \rule{2.2cm}{0.3pt}
  \end{tcolorbox}}

% Bandeau de titre : #1 titre, #2 semaine
\newcommand{\titreblock}[2]{%
\begin{tcolorbox}[colback=teal!10,colframe=teal!70!black,arc=2mm,boxrule=1pt]
\centering
{\LARGE\bfseries #1}\\[2mm]
{\large Robotique -- \gradelabel{} -- Plan de cours de 2 heures}\\[1mm]
Enseignant : Francis Ki \quad|\quad #2
\end{tcolorbox}}

% Tableau du minutage : 6 contenus (accroche, mini-leçon, guidée, autonome, mise en commun, exit ticket)
\newcommand{\minutage}[6]{%
\begin{tabularx}{\linewidth}{>{\bfseries}p{2.2cm}p{2.6cm}Y p{3cm}}
\toprule
Minutes & Étape & Contenu & Organisation\\
\midrule
0 -- 10 & Accroche & #1 & Classe entière\\
10 -- 30 & Mini-leçon & #2 & Classe entière\\
30 -- 60 & Pratique guidée & #3 & Collectif puis binômes\\
60 -- 95 & Pratique autonome & #4 & Individuel sur poste\\
95 -- 110 & Mise en commun & #5 & Classe entière\\
110 -- 120 & Exit ticket & #6 & Individuel sur poste\\
\bottomrule
\end{tabularx}}

% Étayage progressif : #1 aide 1, #2 aide 2, #3 aide 3
\newcommand{\aides}[3]{%
\textbf{Étayage progressif de l'aide.}
\begin{center}
\begin{tabularx}{\linewidth}{>{\bfseries}p{2.4cm}Y p{4.4cm}}
\toprule
Niveau d'aide & Contenu & Qui intervient\\
\midrule
Aide 1 & #1 & Le binôme seul\\
Aide 2 & #2 & Le binôme avec la carte d'aide\\
Aide 3 & #3 & L'enseignant\\
\bottomrule
\end{tabularx}
\end{center}}

% Les trois groupes de la pratique autonome : #1 tâche A, #2 tâche B, #3 tâche C
\newcommand{\groupes}[3]{%
\textbf{Consigne commune\,:} chaque élève ouvre le fichier de son groupe, le renomme avec son prénom et travaille seul sur son poste. L'enseignant passe d'abord auprès du groupe A pendant 10 min.

\medskip
\begin{tabularx}{\linewidth}{>{\bfseries}p{3cm}Y p{2.2cm}}
\toprule
Groupe & Tâche sur poste & Annexe\\
\midrule
A : en dessous du niveau & #1 & Annexe C\\
B : au niveau & #2 & Annexe D\\
C : au-dessus du niveau & #3 & Annexe E\\
\bottomrule
\end{tabularx}}

% Mise en commun : #1 groupe A, #2 groupe B, #3 groupe C, #4 conclusion
\newcommand{\miseencommun}[4]{%
\begin{etape}{teal}{5. Mise en commun}{95 -- 110 min}
\begin{enumerate}
 \item (2 min) Retour au calme, écrans mis en veille.
 \item (9 min) Trois passages courts de 3 min, avec valorisation :
 \begin{itemize}
  \item \textbf{Groupe A} : #1
  \item \textbf{Groupe B} : #2
  \item \textbf{Groupe C} : #3
 \end{itemize}
 \item (4 min) Conclusion : #4
\end{enumerate}
\end{etape}}

% Clôture : #1 question 1, #2 question 2
\newcommand{\cloture}[2]{%
\begin{etape}{gray}{6. Clôture : exit ticket}{110 -- 120 min}
Formulaire en ligne ou papier (annexe G) :
\begin{enumerate}
 \item #1
 \item #2
 \item Auto-évaluation : \emph{Je maîtrise / Je progresse / J'ai besoin d'aide.}
\end{enumerate}
Lire les réponses avant la séance suivante pour ajuster les groupes de besoin.
\end{etape}}

% Grille d'observation : #1 #2 #3 intitulés de colonnes
\newcommand{\grille}[3]{%
\subsection*{Grille d'observation (enseignant)}
\begin{tabularx}{\linewidth}{Y p{2.4cm} p{2.4cm} p{2.4cm} Y}
\toprule
Élève & #1 & #2 & #3 & Groupe suivant\\
\midrule
 & & & & \\ & & & & \\ & & & & \\ & & & & \\ & & & & \\
\bottomrule
\end{tabularx}}

% Signature de l'exit ticket
\newcommand{\autoeval}{Je me sens : \cb Je maîtrise \quad \cb Je progresse \quad \cb J'ai besoin d'aide}
.
\usepackage[T1]{fontenc}
\usepackage[utf8]{inputenc}
\usepackage{lmodern}
\usepackage[french]{babel}
\usepackage[margin=2cm,headheight=14pt]{geometry}
\usepackage[most]{tcolorbox}
\usepackage{xcolor,tabularx,booktabs,array,enumitem,fancyhdr,amssymb,tikz}
\usetikzlibrary{arrows.meta,positioning}

\pagestyle{fancy}
\fancyhf{}
\lhead{\small Robotique -- \gradelabel{} -- Novembre 2026}
\rhead{\small \leconhdr}
\cfoot{\thepage}
\setlist{itemsep=1pt,topsep=2pt,leftmargin=*}
\renewcommand{\arraystretch}{1.35}
\newcolumntype{Y}{>{\raggedright\arraybackslash}X}
\newcommand{\cb}{$\square$\,}
\newcommand{\fl}{$\rightarrow$}
\newcommand{\lig}{\noindent\rule{\linewidth}{0.3pt}\\[5mm]}

\newtcolorbox{etape}[3]{enhanced,breakable,colback=#1!6,colframe=#1!70!black,
  coltitle=white,fonttitle=\bfseries,title={#2\hfill #3},boxrule=0.8pt,arc=2mm,
  before skip=8pt,after skip=8pt}
\newtcolorbox{astuce}[1][Différenciation et besoins particuliers]{breakable,
  colback=yellow!12,colframe=yellow!50!black,fonttitle=\bfseries\small,
  title={#1},boxrule=0.5pt,arc=1mm,before skip=6pt,after skip=6pt}
\newtcolorbox{code}[1][Code]{breakable,colback=black!4,colframe=black!60,
  fonttitle=\bfseries\small,title={#1},boxrule=0.5pt,arc=1mm,fontupper=\ttfamily\small,
  before skip=6pt,after skip=6pt}
% Blocs de code coloré (écrire le code sans accents : comme dans l'éditeur)
\newtcblisting{arduino}[1][Code Arduino]{listing only,breakable,colback=black!4,
  colframe=black!60,fonttitle=\bfseries\small,title={#1},boxrule=0.5pt,arc=1mm,
  before skip=6pt,after skip=6pt,
  listing options={language=C++,basicstyle=\ttfamily\small,columns=fullflexible,
  keepspaces=true,showstringspaces=false,commentstyle=\color{green!40!black}\itshape,
  keywordstyle=\color{blue!70!black}\bfseries}}
\newtcblisting{python}[1][Code Python]{listing only,breakable,colback=black!4,
  colframe=black!60,fonttitle=\bfseries\small,title={#1},boxrule=0.5pt,arc=1mm,
  before skip=6pt,after skip=6pt,
  listing options={language=Python,basicstyle=\ttfamily\small,columns=fullflexible,
  keepspaces=true,showstringspaces=false,commentstyle=\color{green!40!black}\itshape,
  keywordstyle=\color{blue!70!black}\bfseries}}
\newcommand{\entete}[2]{%
  \begin{tcolorbox}[colback=#1!10,colframe=#1!70!black,boxrule=0.8pt,arc=2mm]
  {\Large\bfseries #2}\hfill Nom : \rule{4cm}{0.3pt}\quad Date : \rule{2.2cm}{0.3pt}
  \end{tcolorbox}}

% Bandeau de titre : #1 titre, #2 semaine
\newcommand{\titreblock}[2]{%
\begin{tcolorbox}[colback=teal!10,colframe=teal!70!black,arc=2mm,boxrule=1pt]
\centering
{\LARGE\bfseries #1}\\[2mm]
{\large Robotique -- \gradelabel{} -- Plan de cours de 2 heures}\\[1mm]
Enseignant : Francis Ki \quad|\quad #2
\end{tcolorbox}}

% Tableau du minutage : 6 contenus (accroche, mini-leçon, guidée, autonome, mise en commun, exit ticket)
\newcommand{\minutage}[6]{%
\begin{tabularx}{\linewidth}{>{\bfseries}p{2.2cm}p{2.6cm}Y p{3cm}}
\toprule
Minutes & Étape & Contenu & Organisation\\
\midrule
0 -- 10 & Accroche & #1 & Classe entière\\
10 -- 30 & Mini-leçon & #2 & Classe entière\\
30 -- 60 & Pratique guidée & #3 & Collectif puis binômes\\
60 -- 95 & Pratique autonome & #4 & Individuel sur poste\\
95 -- 110 & Mise en commun & #5 & Classe entière\\
110 -- 120 & Exit ticket & #6 & Individuel sur poste\\
\bottomrule
\end{tabularx}}

% Étayage progressif : #1 aide 1, #2 aide 2, #3 aide 3
\newcommand{\aides}[3]{%
\textbf{Étayage progressif de l'aide.}
\begin{center}
\begin{tabularx}{\linewidth}{>{\bfseries}p{2.4cm}Y p{4.4cm}}
\toprule
Niveau d'aide & Contenu & Qui intervient\\
\midrule
Aide 1 & #1 & Le binôme seul\\
Aide 2 & #2 & Le binôme avec la carte d'aide\\
Aide 3 & #3 & L'enseignant\\
\bottomrule
\end{tabularx}
\end{center}}

% Les trois groupes de la pratique autonome : #1 tâche A, #2 tâche B, #3 tâche C
\newcommand{\groupes}[3]{%
\textbf{Consigne commune\,:} chaque élève ouvre le fichier de son groupe, le renomme avec son prénom et travaille seul sur son poste. L'enseignant passe d'abord auprès du groupe A pendant 10 min.

\medskip
\begin{tabularx}{\linewidth}{>{\bfseries}p{3cm}Y p{2.2cm}}
\toprule
Groupe & Tâche sur poste & Annexe\\
\midrule
A : en dessous du niveau & #1 & Annexe C\\
B : au niveau & #2 & Annexe D\\
C : au-dessus du niveau & #3 & Annexe E\\
\bottomrule
\end{tabularx}}

% Mise en commun : #1 groupe A, #2 groupe B, #3 groupe C, #4 conclusion
\newcommand{\miseencommun}[4]{%
\begin{etape}{teal}{5. Mise en commun}{95 -- 110 min}
\begin{enumerate}
 \item (2 min) Retour au calme, écrans mis en veille.
 \item (9 min) Trois passages courts de 3 min, avec valorisation :
 \begin{itemize}
  \item \textbf{Groupe A} : #1
  \item \textbf{Groupe B} : #2
  \item \textbf{Groupe C} : #3
 \end{itemize}
 \item (4 min) Conclusion : #4
\end{enumerate}
\end{etape}}

% Clôture : #1 question 1, #2 question 2
\newcommand{\cloture}[2]{%
\begin{etape}{gray}{6. Clôture : exit ticket}{110 -- 120 min}
Formulaire en ligne ou papier (annexe G) :
\begin{enumerate}
 \item #1
 \item #2
 \item Auto-évaluation : \emph{Je maîtrise / Je progresse / J'ai besoin d'aide.}
\end{enumerate}
Lire les réponses avant la séance suivante pour ajuster les groupes de besoin.
\end{etape}}

% Grille d'observation : #1 #2 #3 intitulés de colonnes
\newcommand{\grille}[3]{%
\subsection*{Grille d'observation (enseignant)}
\begin{tabularx}{\linewidth}{Y p{2.4cm} p{2.4cm} p{2.4cm} Y}
\toprule
Élève & #1 & #2 & #3 & Groupe suivant\\
\midrule
 & & & & \\ & & & & \\ & & & & \\ & & & & \\ & & & & \\
\bottomrule
\end{tabularx}}

% Signature de l'exit ticket
\newcommand{\autoeval}{Je me sens : \cb Je maîtrise \quad \cb Je progresse \quad \cb J'ai besoin d'aide}
.
\usepackage[T1]{fontenc}
\usepackage[utf8]{inputenc}
\usepackage{lmodern}
\usepackage[french]{babel}
\usepackage[margin=2cm,headheight=14pt]{geometry}
\usepackage[most]{tcolorbox}
\usepackage{xcolor,tabularx,booktabs,array,enumitem,fancyhdr,amssymb,tikz}
\usetikzlibrary{arrows.meta,positioning}

\pagestyle{fancy}
\fancyhf{}
\lhead{\small Robotique -- \gradelabel{} -- Novembre 2026}
\rhead{\small \leconhdr}
\cfoot{\thepage}
\setlist{itemsep=1pt,topsep=2pt,leftmargin=*}
\renewcommand{\arraystretch}{1.35}
\newcolumntype{Y}{>{\raggedright\arraybackslash}X}
\newcommand{\cb}{$\square$\,}
\newcommand{\fl}{$\rightarrow$}
\newcommand{\lig}{\noindent\rule{\linewidth}{0.3pt}\\[5mm]}

\newtcolorbox{etape}[3]{enhanced,breakable,colback=#1!6,colframe=#1!70!black,
  coltitle=white,fonttitle=\bfseries,title={#2\hfill #3},boxrule=0.8pt,arc=2mm,
  before skip=8pt,after skip=8pt}
\newtcolorbox{astuce}[1][Différenciation et besoins particuliers]{breakable,
  colback=yellow!12,colframe=yellow!50!black,fonttitle=\bfseries\small,
  title={#1},boxrule=0.5pt,arc=1mm,before skip=6pt,after skip=6pt}
\newtcolorbox{code}[1][Code]{breakable,colback=black!4,colframe=black!60,
  fonttitle=\bfseries\small,title={#1},boxrule=0.5pt,arc=1mm,fontupper=\ttfamily\small,
  before skip=6pt,after skip=6pt}
% Blocs de code coloré (écrire le code sans accents : comme dans l'éditeur)
\newtcblisting{arduino}[1][Code Arduino]{listing only,breakable,colback=black!4,
  colframe=black!60,fonttitle=\bfseries\small,title={#1},boxrule=0.5pt,arc=1mm,
  before skip=6pt,after skip=6pt,
  listing options={language=C++,basicstyle=\ttfamily\small,columns=fullflexible,
  keepspaces=true,showstringspaces=false,commentstyle=\color{green!40!black}\itshape,
  keywordstyle=\color{blue!70!black}\bfseries}}
\newtcblisting{python}[1][Code Python]{listing only,breakable,colback=black!4,
  colframe=black!60,fonttitle=\bfseries\small,title={#1},boxrule=0.5pt,arc=1mm,
  before skip=6pt,after skip=6pt,
  listing options={language=Python,basicstyle=\ttfamily\small,columns=fullflexible,
  keepspaces=true,showstringspaces=false,commentstyle=\color{green!40!black}\itshape,
  keywordstyle=\color{blue!70!black}\bfseries}}
\newcommand{\entete}[2]{%
  \begin{tcolorbox}[colback=#1!10,colframe=#1!70!black,boxrule=0.8pt,arc=2mm]
  {\Large\bfseries #2}\hfill Nom : \rule{4cm}{0.3pt}\quad Date : \rule{2.2cm}{0.3pt}
  \end{tcolorbox}}

% Bandeau de titre : #1 titre, #2 semaine
\newcommand{\titreblock}[2]{%
\begin{tcolorbox}[colback=teal!10,colframe=teal!70!black,arc=2mm,boxrule=1pt]
\centering
{\LARGE\bfseries #1}\\[2mm]
{\large Robotique -- \gradelabel{} -- Plan de cours de 2 heures}\\[1mm]
Enseignant : Francis Ki \quad|\quad #2
\end{tcolorbox}}

% Tableau du minutage : 6 contenus (accroche, mini-leçon, guidée, autonome, mise en commun, exit ticket)
\newcommand{\minutage}[6]{%
\begin{tabularx}{\linewidth}{>{\bfseries}p{2.2cm}p{2.6cm}Y p{3cm}}
\toprule
Minutes & Étape & Contenu & Organisation\\
\midrule
0 -- 10 & Accroche & #1 & Classe entière\\
10 -- 30 & Mini-leçon & #2 & Classe entière\\
30 -- 60 & Pratique guidée & #3 & Collectif puis binômes\\
60 -- 95 & Pratique autonome & #4 & Individuel sur poste\\
95 -- 110 & Mise en commun & #5 & Classe entière\\
110 -- 120 & Exit ticket & #6 & Individuel sur poste\\
\bottomrule
\end{tabularx}}

% Étayage progressif : #1 aide 1, #2 aide 2, #3 aide 3
\newcommand{\aides}[3]{%
\textbf{Étayage progressif de l'aide.}
\begin{center}
\begin{tabularx}{\linewidth}{>{\bfseries}p{2.4cm}Y p{4.4cm}}
\toprule
Niveau d'aide & Contenu & Qui intervient\\
\midrule
Aide 1 & #1 & Le binôme seul\\
Aide 2 & #2 & Le binôme avec la carte d'aide\\
Aide 3 & #3 & L'enseignant\\
\bottomrule
\end{tabularx}
\end{center}}

% Les trois groupes de la pratique autonome : #1 tâche A, #2 tâche B, #3 tâche C
\newcommand{\groupes}[3]{%
\textbf{Consigne commune\,:} chaque élève ouvre le fichier de son groupe, le renomme avec son prénom et travaille seul sur son poste. L'enseignant passe d'abord auprès du groupe A pendant 10 min.

\medskip
\begin{tabularx}{\linewidth}{>{\bfseries}p{3cm}Y p{2.2cm}}
\toprule
Groupe & Tâche sur poste & Annexe\\
\midrule
A : en dessous du niveau & #1 & Annexe C\\
B : au niveau & #2 & Annexe D\\
C : au-dessus du niveau & #3 & Annexe E\\
\bottomrule
\end{tabularx}}

% Mise en commun : #1 groupe A, #2 groupe B, #3 groupe C, #4 conclusion
\newcommand{\miseencommun}[4]{%
\begin{etape}{teal}{5. Mise en commun}{95 -- 110 min}
\begin{enumerate}
 \item (2 min) Retour au calme, écrans mis en veille.
 \item (9 min) Trois passages courts de 3 min, avec valorisation :
 \begin{itemize}
  \item \textbf{Groupe A} : #1
  \item \textbf{Groupe B} : #2
  \item \textbf{Groupe C} : #3
 \end{itemize}
 \item (4 min) Conclusion : #4
\end{enumerate}
\end{etape}}

% Clôture : #1 question 1, #2 question 2
\newcommand{\cloture}[2]{%
\begin{etape}{gray}{6. Clôture : exit ticket}{110 -- 120 min}
Formulaire en ligne ou papier (annexe G) :
\begin{enumerate}
 \item #1
 \item #2
 \item Auto-évaluation : \emph{Je maîtrise / Je progresse / J'ai besoin d'aide.}
\end{enumerate}
Lire les réponses avant la séance suivante pour ajuster les groupes de besoin.
\end{etape}}

% Grille d'observation : #1 #2 #3 intitulés de colonnes
\newcommand{\grille}[3]{%
\subsection*{Grille d'observation (enseignant)}
\begin{tabularx}{\linewidth}{Y p{2.4cm} p{2.4cm} p{2.4cm} Y}
\toprule
Élève & #1 & #2 & #3 & Groupe suivant\\
\midrule
 & & & & \\ & & & & \\ & & & & \\ & & & & \\ & & & & \\
\bottomrule
\end{tabularx}}

% Signature de l'exit ticket
\newcommand{\autoeval}{Je me sens : \cb Je maîtrise \quad \cb Je progresse \quad \cb J'ai besoin d'aide}


\begin{document}
\titreblock{Architecture d'un objet connecté}{Semaine du 2 au 6 novembre 2026}

%==================================================================
\section{Informations générales}
\begin{tabularx}{\linewidth}{>{\bfseries}p{3.6cm}Y}
\toprule
Titre & Architecture d'un objet connecté\\
Classe / effectif & Grade 12, 11 à 20 élèves, niveaux hétérogènes, élèves à besoins particuliers\\
Durée & 2 heures (120 min)\\
Lieu & Salle informatique, un poste par élève\\
Matériel et logiciels &
\begin{itemize}
 \item Un poste par élève avec navigateur web.
 \item \textbf{diagrams.net} (draw.io, gratuit, sans compte) ou Google Slides pour les schémas d'architecture.
 \item \textbf{Wokwi} (wokwi.com, gratuit) : un projet ESP32 modèle, en lecture seule, avec un capteur de température, une LED et le moniteur série, pour observer ce que fait l'«~objet~».
 \item Vidéoprojecteur ou tableau interactif.
\end{itemize}\\
Pré-requis &
G10 et G11 : ESP32, entrées et sorties, chaîne d'information (acquérir, traiter, communiquer), notions de réseau (adresse IP, Wi-Fi).\\
\bottomrule
\end{tabularx}

\begin{astuce}[Préparation avant le cours (environ 45 min)]
\begin{enumerate}
 \item Préparer un schéma d'architecture modèle dans draw.io (objet, réseau local, passerelle, internet, serveur, application).
 \item Préparer le projet Wokwi (ESP32 + capteur de température + LED) et vérifier le lien partagé.
 \item Sélectionner trois objets connectés pour les trois niveaux (thermostat connecté, ampoule connectée, serrure connectée).
 \item Préparer les fichiers des annexes C à E et le formulaire de sortie.
\end{enumerate}
\end{astuce}

%==================================================================
\section{Objectifs d'apprentissage}
\begin{tcolorbox}[colback=gray!8,colframe=gray!60!black,arc=1mm,boxrule=0.5pt]
\textbf{Objectifs du programme}
\begin{itemize}
 \item Identifier les fonctions d'un objet connecté.
 \item Comprendre les échanges d'informations.
\end{itemize}
\textbf{À la fin de la séance, l'élève est capable\,:}
\begin{itemize}
 \item de décomposer un objet connecté en blocs fonctionnels (capter, traiter, communiquer, alimenter, interagir)\,;
 \item de représenter l'architecture complète (objet, réseau, serveur, application) avec le sens des flux de données et de commandes\,;
 \item d'argumenter sur les contraintes d'un objet connecté (énergie, latence, sécurité, vie privée).
\end{itemize}
\textbf{Compétences numériques\,:} schématiser une architecture, lire un code court, évaluer les risques liés aux données.
\end{tcolorbox}

%==================================================================
\section{Déroulé minuté (120 min)}
\minutage{Que sait ton thermostat de toi ?}{Fonctions, architecture en couches, flux et contraintes}{Un thermostat connecté en collectif, puis une ampoule connectée en binômes}{Trois niveaux de besoin}{Valoriser chaque groupe}{Trois questions sur poste}

%------------------------------------------------------------------
\begin{etape}{orange}{1. Accroche : que sait ton thermostat de toi ?}{0 -- 10 min}
\begin{enumerate}
 \item (5 min) Présenter un thermostat connecté et sa notification sur téléphone. \emph{Quelles informations circulent ? Qui peut les voir ? Où sont-elles enregistrées ?}
 \item (5 min) Recueillir les idées en nuage de mots. Question-problème : \textbf{«~Que contient un objet connecté et avec qui échange-t-il ?~»}
\end{enumerate}
\end{etape}

%------------------------------------------------------------------
\begin{etape}{blue}{2. Mini-leçon : l'architecture d'un objet connecté}{10 -- 30 min}
\textbf{Support\,:} annexe A. Quatre temps de 5 min.

\textbf{Temps 1 : cinq fonctions dans l'objet.}
\begin{center}
\begin{tabularx}{\linewidth}{>{\bfseries}p{3cm}Y Y}
\toprule
Fonction & Rôle & Exemple (thermostat)\\
\midrule
Alimenter & Fournir l'énergie & Secteur ou piles\\
Capter & Mesurer le monde & Sonde de température\\
Traiter & Décider selon un programme & Microcontrôleur ESP32\\
Communiquer & Échanger avec l'extérieur & Module Wi-Fi\\
Interagir & Afficher ou agir & Écran, relais de chaudière\\
\bottomrule
\end{tabularx}
\end{center}

\textbf{Temps 2 : l'architecture complète.}
\begin{center}
\begin{tikzpicture}[box/.style={draw,rounded corners,thick,minimum height=1.1cm,text width=2.3cm,align=center,font=\small},>={Stealth[length=3mm]},node distance=0.6cm]
\node[box,fill=orange!20] (a) {\textbf{Objet}\\ capteurs, ESP32};
\node[box,fill=blue!15,right=of a] (b) {\textbf{Réseau local}\\ Wi-Fi, box};
\node[box,fill=green!15,right=of b] (c) {\textbf{Internet}\\ serveur, nuage};
\node[box,fill=violet!15,right=of c] (d) {\textbf{Application}\\ téléphone};
\draw[<->,thick] (a)--(b); \draw[<->,thick] (b)--(c); \draw[<->,thick] (c)--(d);
\end{tikzpicture}
\end{center}
Les flèches vont dans les deux sens : les \emph{données} montent (température), les \emph{commandes} descendent (consigne de chauffage).

\textbf{Temps 3 : stocker et traiter où ?} Dans l'objet (rapide, confidentiel, limité) ou dans le nuage (puissant, accessible, dépendant du réseau).

\textbf{Temps 4 : les contraintes.}
\begin{center}
\begin{tabularx}{\linewidth}{>{\bfseries}p{3cm}Y}
\toprule
Contrainte & Question à se poser\\
\midrule
Énergie & Combien l'objet consomme-t-il ? Sur piles, combien de temps ?\\
Latence & Combien de temps entre l'ordre et l'action ?\\
Sécurité & Qui peut commander l'objet ? Les échanges sont-ils chiffrés ?\\
Vie privée & Quelles données personnelles sont collectées ?\\
\bottomrule
\end{tabularx}
\end{center}

\begin{astuce}
Utiliser une couleur par bloc (objet, réseau, nuage, application) dans tous les schémas. Fournir un schéma pré-rempli à trous pour les élèves qui ont besoin d'un appui visuel fort. Autoriser une explication orale enregistrée.
\end{astuce}
\end{etape}

%------------------------------------------------------------------
\begin{etape}{green}{3. Pratique guidée}{30 -- 60 min}
\textbf{Fiche\,:} annexe B. Les élèves travaillent en binômes devant leur poste.

\textbf{3a. Un thermostat connecté en collectif (12 min).} Compléter ensemble le tableau des cinq fonctions, puis dessiner dans draw.io l'architecture complète (projeter le fichier de l'enseignant). Ouvrir le projet Wokwi pour voir ce que fait l'«~objet~» (lecture de température et affichage).

\textbf{3b. Une ampoule connectée en binômes (18 min).} Réaliser le même travail sur un nouvel objet : tableau des fonctions, schéma d'architecture, et sens des flux (données, commandes). Alternance des rôles.

\aides{Le tableau du thermostat et le schéma modèle}{Questions guides : \emph{Qu'est-ce qui mesure ? Qui décide ? Par où passent les données ? Qui envoie une commande ?}}{L'enseignant remplit une ligne avec le binôme, qui poursuit}

\begin{astuce}
Accepter un croquis à la main photographié si le dessin sur poste ralentit l'élève. Toujours commencer par la fonction globale de l'objet.
\end{astuce}
\end{etape}

%------------------------------------------------------------------
\begin{etape}{violet}{4. Pratique autonome différenciée}{60 -- 95 min}
\groupes{Serrure connectée : tableau des cinq fonctions à compléter avec une banque de mots (batterie, lecteur, ESP32, Wi-Fi, gâche) et schéma d'architecture à légender avec des étiquettes mobiles.}{Station météo connectée : tableau des fonctions, schéma d'architecture complet avec le sens des flux (données montantes, commandes descendantes) et trois scénarios d'usage.}{Analyse critique d'un objet connecté au choix : architecture, tableau des quatre contraintes (énergie, latence, sécurité, vie privée), comparaison «~traitement local ou dans le nuage~» et recommandation argumentée.}

\medskip
\textbf{Plan B si un élève termine tôt\,:} chercher un exemple réel de faille de sécurité d'un objet connecté et résumer ce qu'on peut en apprendre (en citant la source).
\end{etape}

%------------------------------------------------------------------
\miseencommun{un élève présente son tableau complété. Valoriser la précision du vocabulaire.}{un binôme projette son schéma et suit le chemin d'une donnée puis d'une commande. Valoriser la clarté.}{un élève présente son analyse des contraintes et sa recommandation. Valoriser l'esprit critique.}{\textit{un objet connecté combine capter, traiter, communiquer ; il s'intègre dans une architecture avec un réseau, un serveur et une application, et pose des questions de sécurité et de vie privée.}}

\cloture{Cite les cinq fonctions d'un objet connecté.}{Pourquoi dit-on que les flèches de l'architecture vont dans les deux sens ?}

%==================================================================
% ANNEXES
%==================================================================
\newpage
\section*{Annexe A : Support résumé de la mini-leçon}
\entete{blue}{Architecture d'un objet connecté}

\begin{tcolorbox}[colback=yellow!10,colframe=yellow!50!black]
Un objet connecté = \textbf{alimenter, capter, traiter, communiquer, interagir}. Il échange avec un \textbf{réseau}, un \textbf{serveur} et une \textbf{application}.
\end{tcolorbox}

\begin{center}
\begin{tikzpicture}[box/.style={draw,rounded corners,thick,minimum height=1.1cm,text width=2.3cm,align=center,font=\small},>={Stealth[length=3mm]},node distance=0.6cm]
\node[box,fill=orange!20] (a) {Objet};
\node[box,fill=blue!15,right=of a] (b) {Réseau local};
\node[box,fill=green!15,right=of b] (c) {Internet, serveur};
\node[box,fill=violet!15,right=of c] (d) {Application};
\draw[<->,thick] (a)--(b); \draw[<->,thick] (b)--(c); \draw[<->,thick] (c)--(d);
\end{tikzpicture}
\end{center}

\begin{tabularx}{\linewidth}{>{\bfseries}p{3cm}Y}
\toprule
Contrainte & Question\\
\midrule
Énergie & Quelle autonomie ?\\
Latence & Quel délai entre ordre et action ?\\
Sécurité & Qui peut commander l'objet ?\\
Vie privée & Quelles données personnelles ?\\
\bottomrule
\end{tabularx}

%==================================================================
\newpage
\section*{Annexe B : Fiche de pratique guidée}
\entete{green}{Pratique guidée : décrire un objet connecté}

\subsection*{3a. Le thermostat connecté (collectif)}
\begin{tabularx}{\linewidth}{>{\bfseries}p{3cm}Y Y}
\toprule
Fonction & Composant du thermostat & Information échangée\\
\midrule
Alimenter & & \\ Capter & & \\ Traiter & & \\ Communiquer & & \\ Interagir & & \\
\bottomrule
\end{tabularx}

\subsection*{3b. L'ampoule connectée (binôme)}
\begin{tabularx}{\linewidth}{>{\bfseries}p{3cm}Y Y}
\toprule
Fonction & Composant de l'ampoule & Donnée ou commande\\
\midrule
Alimenter & & \\ Capter & & \\ Traiter & & \\ Communiquer & & \\ Interagir & & \\
\bottomrule
\end{tabularx}

Sur une diapositive, dessine l'architecture et marque en couleur les données (montantes) et les commandes (descendantes).

\begin{tcolorbox}[colback=yellow!10,colframe=yellow!50!black,title={Cartes d'aide à utiliser dans l'ordre},fonttitle=\bfseries\small]
\textbf{Aide 1 :} regarde le tableau du thermostat et le schéma de l'annexe A.\\
\textbf{Aide 2 :} Qu'est-ce qui mesure ? Qui décide ? Par où passent les données ? Qui envoie une commande ?\\
\textbf{Aide 3 :} demande à l'enseignant de remplir une ligne avec toi.
\end{tcolorbox}

%==================================================================
\newpage
\section*{Annexe C : Niveau A (en dessous du niveau)}
\entete{orange}{Niveau A : la serrure connectée}

\textbf{Banque :} batterie, lecteur de badge, ESP32, Wi-Fi, gâche électrique.

\begin{tabularx}{\linewidth}{>{\bfseries}p{3cm}Y}
\toprule
Fonction & Composant de la serrure\\
\midrule
Alimenter & \\ Capter & \\ Traiter & \\ Communiquer & \\ Interagir & \\
\bottomrule
\end{tabularx}

\subsection*{Légende le schéma (glisse les étiquettes)}
Serrure \fl{} \dotfill{} \fl{} \dotfill{} \fl{} \dotfill

\subsection*{Vrai ou faux ?}
1. Une serrure connectée peut recevoir une commande depuis un téléphone. \cb Vrai \cb Faux\\
2. Les flèches du schéma ne vont que dans un sens. \cb Vrai \cb Faux

%==================================================================
\newpage
\section*{Annexe D : Niveau B (au niveau)}
\entete{blue}{Niveau B : la station météo connectée}

\begin{tabularx}{\linewidth}{>{\bfseries}p{3cm}Y Y}
\toprule
Fonction & Composant & Information\\
\midrule
Alimenter & & \\ Capter & & \\ Traiter & & \\ Communiquer & & \\ Interagir & & \\
\bottomrule
\end{tabularx}

\subsection*{Trois scénarios}
\begin{tabularx}{\linewidth}{>{\bfseries}p{4cm}Y Y}
\toprule
Scénario & Donnée ou commande & Chemin parcouru\\
\midrule
L'utilisateur consulte la température & & \\
La station envoie une mesure toutes les 10 min & & \\
L'utilisateur change la fréquence des mesures & & \\
\bottomrule
\end{tabularx}

Dessine le schéma complet sur une diapositive.

%==================================================================
\newpage
\section*{Annexe E : Niveau C (au-dessus du niveau)}
\entete{violet}{Niveau C : analyse critique}

\textbf{Mission :} choisis un objet connecté (bracelet, caméra, enceinte, thermostat).

\begin{tabularx}{\linewidth}{>{\bfseries}p{3cm}Y Y}
\toprule
Contrainte & Analyse pour mon objet & Risque ou limite\\
\midrule
Énergie & & \\ Latence & & \\ Sécurité & & \\ Vie privée & & \\
\bottomrule
\end{tabularx}

\subsection*{Local ou nuage ?}
\begin{tabularx}{\linewidth}{>{\bfseries}p{3.4cm}Y Y}
\toprule
Critère & Traitement dans l'objet & Traitement dans le nuage\\
\midrule
Rapidité & & \\ Confidentialité & & \\ Puissance de calcul & & \\ Dépendance au réseau & & \\
\bottomrule
\end{tabularx}

Ma recommandation (cinq lignes maximum) : \lig \lig \lig

\textbf{Défi bonus :} explique ce qui se passe si le réseau tombe en panne.

%==================================================================
\newpage
\section*{Annexe F : Corrections et attendus (enseignant)}

\subsection*{Annexe B}
Thermostat : alimenter (secteur ou piles), capter (sonde), traiter (ESP32), communiquer (Wi-Fi), interagir (écran, relais). Ampoule : alimenter (secteur), capter (parfois rien ou un capteur de présence), traiter (puce), communiquer (Wi-Fi, Bluetooth), interagir (LEDs). Flux : mesures montantes, commandes descendantes.

\subsection*{Annexe C}
Batterie ; lecteur de badge ; ESP32 ; Wi-Fi ; gâche électrique. Schéma : serrure \fl{} réseau \fl{} serveur. Vrai, Faux.

\subsection*{Annexe D}
Capteurs : température, humidité, pression ; ESP32 ; Wi-Fi ; écran. Scénario 1 : application $\to$ serveur $\to$ lecture de la dernière mesure. Scénario 2 : station $\to$ réseau $\to$ serveur. Scénario 3 : application $\to$ serveur $\to$ commande $\to$ station.

\subsection*{Annexe E}
Une bonne analyse relève au moins un risque par contrainte (par exemple batterie qui se vide, données de santé exposées, commande piratée). Local : rapide et confidentiel, nuage : puissant mais dépendant du réseau.\\
\textbf{Critères de réussite :} quatre contraintes traitées, recommandation argumentée.

%==================================================================
\newpage
\section*{Annexe G : Exit ticket et grille d'observation}
\entete{gray}{Exit ticket : architecture}

\begin{enumerate}
 \item Cite les cinq fonctions d'un objet connecté.\\ \lig
 \item Pourquoi dit-on que les flèches de l'architecture vont dans les deux sens ?\\ \lig
 \item \autoeval
\end{enumerate}
\textbf{Attendus :} 1. Alimenter, capter, traiter, communiquer, interagir. 2. Les données montent de l'objet vers l'application, les commandes descendent de l'application vers l'objet.

\bigskip
\grille{Cinq fonctions}{Architecture}{Flux}

\end{document}
